% Generated draft: requires mathematical and source review.
\documentclass[11pt]{article}
\usepackage[margin=25mm]{geometry}
\usepackage{fontspec}
\setmainfont{Latin Modern Roman}
\usepackage{amsmath,amssymb,mathtools,unicode-math}
\setmathfont{Latin Modern Math}
\usepackage{enumitem,needspace,xurl,hyperref}
\setlength{\parindent}{0pt}
\setlength{\parskip}{5pt}
\setlength{\emergencystretch}{3em}
\newcommand{\N}{\mathbb{N}}
\newcommand{\Z}{\mathbb{Z}}
\newcommand{\Q}{\mathbb{Q}}
\newcommand{\R}{\mathbb{R}}
\newcommand{\C}{\mathbb{C}}
\newcommand{\F}{\mathbb{F}}
\newcommand{\A}{\mathbb{A}}
\newcommand{\PP}{\mathbb{P}}
\newcommand{\E}{\mathbb{E}}
\newcommand{\eps}{\varepsilon}
\newcommand{\dd}{\,\mathrm{d}}
\begin{document}
\section*{UnsolvedMath ID: 800216}
\textbf{Source problem code:} catalogue entry
\subsection{Definitions and mathematical statement}
An open subset of a complex space is equipped with its inherited complex-space structure. An increasing sequence \((U_n)_{n\geq 1}\) means that \(U_n\subseteq U_{n+1}\) for every \(n\geq 1\), with equality allowed.

Let \(X\) be a Stein complex space, and let \((U_n)_{n\geq 1}\) be an increasing sequence of open subsets of \(X\), each of which is Stein. Set
\[
U=\bigcup_{n=1}^{\infty}U_n.
\]
Must \(U\), with its inherited complex-space structure, be Stein for every such \(X\) and sequence \((U_n)\)?

The term Stein is used in its standard complex-analytic sense. The source does not specify whether its convention for complex spaces requires reducedness; this convention remains unresolved, and the question is not restricted here to complex manifolds.

The source additionally reports that, if the ambient space \(X\) is instead a complex manifold not assumed to be Stein, the analogous assertion is false in general. This is a supplied background assertion, not an independently verified result.
\subsection{Short English statement}
Must every increasing countable union of Stein open subsets of a Stein complex space be Stein? The source reports a negative answer in general when the ambient space is merely a complex manifold. Its convention concerning reduced complex spaces is unspecified.
\subsection{Sources}
\begin{itemize}
\item Local UnsolvedMath registry: permanent ID \(800216\), exact source code \texttt{catalogue entry}, title ``Several Complex Variables Problem List --- Problem 12.'' The supplied record labels the problem open; that label is supplied curation, not an independently verified current status.
\item Buzzard, Fornaess, and Krantz, \emph{The Several Complex Variables Problem List} (1995), Problem 12. The registry attributes the problem to these authors and gives proposed year 1995. Supplied source URL: \url{https://arxiv.org/abs/math/9509202}.
\item The supplied provenance records access on 2026-10-16 and extraction by \texttt{arxiv\_latex}. These are source-provided provenance details; no independent source visit or verification is claimed.
\end{itemize}
\end{document}
